% sec-11: abbreviations, availability, contributions, citation, references.
\section*{Back Matter}
\phantomsection\addcontentsline{toc}{section}{Back Matter}

\bmhead{Abbreviations}

\begin{apptable}
\begin{tabularx}{\textwidth}{>{\raggedright\arraybackslash}p{1.5cm} >{\raggedright\arraybackslash}X >{\raggedright\arraybackslash}p{1.5cm} >{\raggedright\arraybackslash}X}
\headrow \thead{Short} & \thead{Expansion} & \thead{Short} & \thead{Expansion}\\
\midrule
AE & Adverse event & MTD & Maximum tolerated dose \\
AMP & Accelerating Medicines Partnership & NCI & National Cancer Institute \\
ARPA-H & Advanced Research Projects Agency for Health & NIH & National Institutes of Health \\
CTCAE & Common Terminology Criteria for Adverse Events & NSF & National Science Foundation \\
CTEP & Cancer Therapy Evaluation Program & ODE & Ordinary differential equation \\
ctDNA & Circulating tumour DNA & OS & Overall survival \\
DLT & Dose-limiting toxicity & PD & Pharmacodynamics \\
DOI & Digital object identifier & PDAC & Pancreatic ductal adenocarcinoma \\
FNIH & Foundation for the National Institutes of Health & PFS & Progression-free survival \\
FRO & Focused research organization & PK & Pharmacokinetics \\
HHMI & Howard Hughes Medical Institute & QSP & Quantitative systems pharmacology \\
IIT & Investigator-initiated trial & RP2D & Recommended Phase 2 dose \\
IND & Investigational new drug & SAE & Serious adverse event \\
IRB & Institutional review board & SBIR & Small Business Innovation Research \\
ISGPS & International Study Group of Pancreatic Surgery & TIP & Technology, Innovation and Partnerships \\
LLM & Large language model & VVUQ & Verification, validation and uncertainty quantification \\
\end{tabularx}
\end{apptable}
\tabcap{Abbreviations used in this paper, set in two column pairs at the body
measure. Trial-specific terms follow the definitions in the published Phase 1
protocol rather than being redefined here.}

\bmhead{Data and Code Availability}

{\sloppy Every work cited to a digital object identifier is openly deposited.
The repository is at
\url{https://github.com/kevinkawchak/physical-ai-oncology-trials},
version v4.4.0 \cite{repo2026}.\par} It contains the ten application file sets with
their email text and Overleaf bundles, the twenty figure specifications across
five machine-readable diagram platforms, the three build stages of this paper,
the master prompt and the thirteen sub-prompts, and the TikZ source of every
figure printed here. No raster image is used anywhere in this paper or in any of
the ten attachments.

\bmhead{Author Contributions and Conflicts}

Kevin Kawchak is the sole author. He is chief executive of ChemicalQDevice and
holds no equity in Revolution Medicines or in any robotic surgery vendor. No
sponsor, contract research organization, institutional review board, regulator,
or medical society reviewed this paper or authorised any statement in it. The
work was adapted using Claude Code Opus 5.

\bmhead{Citation}

Kawchak K. 10 Funding Applications: A Phase 1, First-in-Human, PDAC Clinical
Trial Protocol of a LLM-Directed Robotic Whipple with Daraxonrasib (RMC-6236).
Zenodo. 2026; \dlink{10.5281/zenodo.xxxxxxxx}.

\vspace{0.2em}
\apprule{0.5pt}
\vspace{-0.1em}

\bmhead{References}
\vspace{-0.35em}
\begin{apprefs}
\bibliographystyle{unsrt}
\bibliography{references}
\end{apprefs}
